% !TeX root = ../induk.tex
\begin{table}[H]
\centering
\caption{Perbandingan Anggota Klaster dengan Kategori IPM (2024)} \label{tab:ipm_cross_2024}
\begin{tabular}{|c|p{2cm}|p{2cm}|p{2cm}|p{2cm}|}
\hline
\textbf{Klaster K-Means} & \textbf{Sangat Tinggi} & \textbf{Tinggi} & \textbf{Sedang} & \textbf{Rendah} \\ \hline
\textbf{Klaster 1} & - & Bogor, Ciamis, Kuningan, Cirebon, Subang, Purwakarta, Bandung Barat & Sukabumi, Cianjur, Garut, Tasikmalaya, Majalengka, Indramayu, Pangandaran, Pandeglang, Lebak & - \\ \hline
\textbf{Klaster 2} & Kota Jakarta Selatan, Kota Jakarta Timur, Kota Jakarta Pusat, Kota Jakarta Barat, Kota Jakarta Utara, Kota Bandung, Kota Bekasi, Kota Depok, Kota Tangerang, Kota Tangerang Selatan & - & - & - \\ \hline
\textbf{Klaster 3} & Kota Cimahi & Kep. Seribu, Bandung, Sumedang, Karawang, Bekasi, Kota Bogor, Kota Sukabumi, Kota Cirebon, Kota Tasikmalaya, Kota Banjar, Tangerang, Kota Cilegon, Kota Serang & Serang & - \\ \hline
\end{tabular}
\end{table}